\chapter{Conclusion and Outlook}
\label{ch:conclusion}

\section{Synthesis of the research findings}
\label{sec:assessment_observed_findings}

This dissertation developed an integrated approach to physically constrained machine learning surrogates of activated sludge systems for process optimization. The central engineering problem was to obtain rapid treatment predictions while retaining a consistent account of the material that passes through biological reactors and solids separation. The research addressed this problem through three connected assessments. Chapter~\ref{ch:projection} compared statistical reactor predictors before and after physical projection. Chapter~\ref{ch:interpretable} developed an explicit regression whose operating, loading, and component relationships could be examined. Chapter~\ref{ch:plant} extended prediction and projection to a connected treatment plant and compared the operating decisions selected by surrogate and direct mechanistic searches.

The component state provided the common foundation. The adopted mechanistic model represented twenty component concentrations and twenty-eight reaction processes. Keeping the component response retained distinctions among organic substrate, nitrogen species, phosphorus storage, biomass, and solids. The fixed reporting map then produced COD, TN, TP, and TSS. This order connected the quantities commonly used to assess treatment with the detailed material accounting used to impose physical constraints. It also allowed the same predicted state to support balance checks, treatment reporting, and operating evaluation.

\subsection{Physical enforcement across different reactor predictors}

The reactor benchmark showed that projection could supply a common physical enforcement procedure across thirteen statistical learners. Every raw assessment state failed the specified mass conservation tolerance. Some predictors returned non-negative concentrations throughout the largest-design assessment, while others produced negative components frequently. Applying mass conservation and non-negativity together returned states that passed both checks in all 766,350 assessed projection instances. The maximum invariant residual was $2.06\times10^{-13}$ for the learning-domain assessment and $1.72\times10^{-13}$ for the exterior assessment. The result established a practical numerical connection between statistical component prediction and the physical accounting of the reactor.

Projection also redistributed prediction error among components. At the largest design size, it improved 188 of the 260 predictor--component comparisons. The aggregate component nRMSE decreased for five predictors and increased for eight. Every predictor had lower TN and TP root error after projection, while COD and TSS root errors increased. These outcomes show why the effect of physical enforcement is most informative when examined at both component and reported-quantity levels. A single aggregate score compresses changes that have different implications for nutrient prediction, biomass, and solids.

The learning curves showed that the available simulation budget influenced predictor choice. XGBoost and LightGBM had the lowest observed raw errors at 500 cases. The multilayer perceptron led at 10,000 cases, with raw and projected component nRMSE values of 0.150 and 0.149. Its projected error increased to 0.251 under mild exterior inputs and 0.521 under severe exterior inputs. All exterior projected states still passed the physical checks. The benchmark therefore provided separate evidence about learning behavior, response to exterior inputs, and constraint enforcement. This separation gives engineers a clear basis for selecting a predictor for a particular data budget and operating domain.

\subsection{Inspectable prediction through structured regression}

ICSOR added an explicit representation of operating dependence and influent loading to the physical prediction framework. Its second-order driver separated baseline, operating, loading, curvature, and interaction terms. Learned component coupling connected these terms to the complete effluent state. Affine and constrained projections then returned the physically checked prediction. This structure made it possible to follow an operating input from its polynomial contribution through component coupling and output enforcement.

The model performed particularly well when fewer simulated cases were available. At 500 cases, its mean held-out pooled composite root error was 10.79, compared with 11.26 for XGBoost and 11.45 for LightGBM. It also had the lowest reported normalized root-error learning-curve area of 6.33. At 10,000 cases, the repeated-split error was 5.94 for ICSOR and 4.62 for the multilayer perceptron. LightGBM had the best mean rank across the broader assessment. The fixed partition gave a similar ordering, with pooled composite root errors of 5.98 for ICSOR and 4.38 for the multilayer perceptron. These results identify an engineering choice between an explicit approximation with favorable low-data behavior and more flexible predictors with lower full-size error.

The deployment assessment identified the role of each physical operation. Affine projection returned non-negative states in 436 of the 2000 fixed-partition cases. The remaining 1564 cases required the constrained linear projection. Every final state passed the adopted numerical balance and non-negativity checks. The complete procedure thus combined structured regression with systematic output enforcement. The branch counts also quantified the additional calculation required to return the final state.

The coefficient analysis showed selective influent curvature in the raw COD mapping together with dense shared component coupling. Forty-four of 210 unique COD influent-curvature entries and 371 of 380 shared off-diagonal coupling entries exceeded their respective reporting thresholds. Named blocks and physical-unit coefficients made these fitted relationships available for inspection. Analytical expressions described the raw and affine sensitivities. The two-component illustration in Chapter~\ref{ch:interpretable} showed how conservation can reverse one raw slope and how an active non-negativity boundary can make the final local slope zero. Interpretation consequently follows the complete prediction calculation and the operating point at which it is examined.

\subsection{Connected prediction and operating optimization}

The connected plant formulation extended physical enforcement from one effluent vector to a joint response containing the mixer, five biological reactors, clarifier overflow, clarifier underflow, and stored solids. Hydraulic connections, stage invariants, flow-weighted separator balances, soluble-component continuity, and inventory bounds linked these outputs. The resulting formulation preserved information about both the liquid treatment path and the recycled solids that influence upstream conditions.

The development and holdout designs retained 7386 and 1845 mechanistic states. On the holdout, joint projection reduced aggregate nRMSE from 0.04690 to 0.04306, a relative improvement of 8.19\%. Aggregate nMAE decreased by 15.65\%. Projection improved 30 of the 32 location--composite nRMSE values. The largest relative composite improvement occurred for TP, and overflow TP error decreased by 39.10\%. The two increases occurred for COD in the final reactor and overflow. These results support the value of a joint plant response while preserving the location-specific information needed to understand its treatment predictions.

Operating assessment introduced the controls selected by each search as an additional object of evaluation. Independent mechanistic calculations supplied the treatment outcomes at those controls. In scenario S4, the surrogate predicted TN and TP concentrations of 2.67 and 2.08 mg L$^{-1}$, while the mechanistic concentrations were 7.47 and 4.68 mg L$^{-1}$. This finding made nutrient prediction at the selected decision directly relevant to the operating comparison. It also showed the practical value of calculating treatment outcomes at the settings that would actually be chosen.

The direct route had lower common-reference objectives in the nominal case and all ten influent scenarios. Nominal objectives were 0.8647 for Extended ICSOR and 0.7862 for Smooth NLP. The direct route generally selected a larger treatment capacity and accepted greater resource contributions to obtain lower quality cost. Scenario S10 illustrated this choice particularly clearly. Its quality cost decreased from 1.4748 to 0.3699, while the weighted resource contribution increased from 0.0557 to 0.4191. The total objective decreased from 0.7931 to 0.6040. In S8, the direct route instead achieved a slightly lower total objective through a lower resource contribution while all four effluent composites were slightly higher. Presenting the objective and its contributions together explained the treatment priorities embodied in each decision.

The connected surrogate shortened the recorded optimization interval. Mean times over the ten scenarios were 1.726 s for Extended ICSOR and 5.236 s for Smooth NLP, giving a ratio of approximately 3.03. This computational result accompanied higher mechanistic objectives for the surrogate decisions. The combined evidence identifies rapid candidate generation and mechanistic decision evaluation as complementary tasks. Their coordinated use is a direct route toward efficient process optimization with explicit treatment and resource priorities.

\section{The integrated methodological contribution}
\label{sec:assessment_conclusions}

The three assessments establish a progression from component prediction to operating choice. The reactor benchmark supplies a common comparison of learners and physical adjustments. The structured regression makes the fitted response and its sensitivities visible. The connected formulation preserves unit relationships and extends the calculation to plant decisions. Projection provides a shared link throughout this progression by converting a raw statistical response into a state that satisfies the specified physical relations.

Table~\ref{tab:integrated_observed_findings} brings together the principal findings and their engineering uses. Each numerical comparison belongs to the assessment in which it was calculated. The reactor errors use training-based component scales, the structured comparison pools four native composite coordinates, and the plant holdout uses location-specific reference ranges. Together, these measures describe different aspects of the same engineering problem.

\begin{table}[htbp]
\centering\small
\caption{Synthesis of the observed findings and their roles in activated sludge process optimization.}
\label{tab:integrated_observed_findings}
\begin{tabular}{>{\raggedright\arraybackslash}p{0.23\textwidth}>{\raggedright\arraybackslash}p{0.39\textwidth}>{\raggedright\arraybackslash}p{0.28\textwidth}}
\toprule
Assessment & Principal finding & Engineering use\\
\midrule
Reactor projection & All 766,350 assessed instances passed both imposed physical checks & Return component states with consistent inventories and non-negative concentrations\\
Reactor prediction & Projection improved 188 of 260 component comparisons and five of thirteen aggregate scores at the largest size & Select learners and assess adjustments for the intended treatment quantities\\
Structured regression & Lowest mean composite error at 500 cases and explicit operating, loading, and coupling terms & Examine fitted treatment relationships with a limited simulation budget\\
Structured deployment & Affine projection sufficed in 21.8\% of fixed-partition cases and constrained projection handled 78.2\% & Follow the complete calculation that supplies the final checked state\\
Connected prediction & Holdout nRMSE decreased by 8.19\%, with improvement in 30 of 32 local scores & Reconcile predictions across biological treatment and solids separation\\
Operating comparison & Mean scenario search times were 1.726 and 5.236 s, while all direct decisions had lower reference objectives & Combine rapid candidate generation with mechanistic treatment evaluation\\
\bottomrule
\end{tabular}
\end{table}

The integration also clarifies how process knowledge enters machine learning. Reaction stoichiometry supplies invariant relations. Flow and equipment boundaries supply transport and inventory relations. Statistical fitting approximates the remaining response over the sampled inputs. Interpretation examines the resulting mapping, and operating evaluation applies the engineering priorities. These contributions work together through defined quantities and checks. Their value lies in making the predicted treatment state and the selected decision open to examination.

For process optimization, the framework supports a sequence of screening, candidate selection, physical adjustment, and reference evaluation. Component-resolved predictions allow nutrient forms, biomass, and solids to be examined alongside the reported effluent quantities. The connected response accounts for recycle effects and solids storage. Independent evaluation then supplies treatment outcomes under the controls selected by the search. This sequence makes the time required to explore operating alternatives useful within a decision process that retains the original treatment objective.

\section{Directions for future research}

\subsection{Adaptive sampling around operating decisions}

Future studies can generate new mechanistic cases where the surrogate is most influential in the operating search. A first surrogate can screen the declared domain and identify candidate controls. Mechanistic calculations at those candidates can then supply additional training states. Sampling can concentrate on large nutrient discrepancies, strong projection displacements, and regions close to treatment requirements. Repeating fitting and operating assessment would connect the simulation budget to the decisions it is intended to support.

An informative comparison would hold the mechanistic evaluation budget fixed and compare uniform sampling, error-based sampling, and decision-based sampling. The outcomes should include effluent errors at selected controls, the common-reference objective, the number of accepted decisions, and the complete time required to deliver them. The S4 nutrient discrepancies provide a concrete target for such a study. Additional cases near those selected operating conditions could test whether focused sampling improves TN and TP prediction where the operating choice depends on them.

\subsection{Projection aligned with treatment priorities}

The reactor results motivate further study of projection weights and the quantities they emphasize. A projection that minimizes changes in component coordinates can redistribute error differently from an assessment based on reported nutrients or solids. Future formulations can compare component scaling, composite-sensitive weighting, and weights based on prediction uncertainty. Each formulation should retain the same mass conservation and non-negativity requirements so that changes in treatment accuracy can be attributed to the adjustment metric.

The connected plant offers an additional test of this idea. Weighting can distinguish internal state accuracy from final-effluent accuracy while preserving the unit connections. Studies can assess whether nutrient-focused weights improve selected-control predictions without producing large errors in biomass, solids inventory, or other treatment locations. This would make the relation between physical adjustment and the intended engineering decision more explicit.

\subsection{Uncertainty-aware prediction and operating safeguards}

Uncertainty can be developed alongside the projected point prediction. Repeated fits, model ensembles, or probabilistic component predictions can describe variation in the estimated treatment response. Each sampled state can be projected onto the same physical set. The resulting distributions can then be evaluated for their coverage of mechanistic reference concentrations, particularly for TN, TP, and overflow TSS.

Operating studies can use these distributions to formulate treatment safeguards. For example, a candidate setting can be selected to keep an upper predictive bound for TP below a specified treatment requirement. Mechanistic and plant-based evaluations can test the frequency with which those safeguards are met. Separating uncertainty associated with surrogate fitting, influent measurements, and mechanistic parameters would also show which information most improves the decision.

\subsection{Stable interpretation of the complete response}

Further interpretation work can examine the stability of coefficient groups across repeated training samples and operating domains. The present structure provides natural groups for operating terms, loading terms, curvature, and coupling. Resampling can show which groups remain prominent and which individual coefficients vary. Comparisons based on the complete predicted response can help identify different coefficient combinations that describe similar treatment behavior.

Sensitivity studies should also follow the final projection. Derivatives can be evaluated at representative operating points and checked against direct perturbations of the returned concentrations. Mapping the active constraints would show where a nutrient response changes smoothly and where a projection boundary alters its slope. Comparing these projected sensitivities with mechanistic sensitivities would connect model inspection more closely to aeration, recycle, and wasting decisions.

\subsection{Dynamic prediction and control}

A dynamic extension can replace steady-state targets with trajectories of component concentrations and stored solids. The governing balances would include the time derivatives of reactor and clarifier inventories. Training data can represent changing influent loads, aeration schedules, recycle settings, and sludge wasting. Physical enforcement can then account for the material transported and transformed during each prediction interval.

This extension would support studies of predictive control under daily influent variation and load disturbances. The evaluation can include transient nutrient peaks, recovery time, inventory changes, oxygen demand, and the computation required at each control update. The connected component representation provides a useful starting point because it already retains the locations and inventories that influence delayed treatment responses.

\subsection{Plant validation, topology, and resource models}

Validation with measured plant data is a major next step. Studies can combine influent fractionation, calibrated mechanistic parameters, unit concentrations, clarifier solids measurements, and recorded controls. A prospective assessment can fix the fitting and operating procedure before collecting the final evaluation period. Testing across seasons and loading conditions would provide evidence about how the component surrogate transfers to observed treatment behavior.

Operating objectives can also be developed for specific facilities. An installed plant has defined reactor volumes, aeration capacity, recycle capacity, and discharge requirements. Electricity demand, chemical use, sludge handling, and operating costs can replace or supplement the engineering proxies used here. Multiobjective studies can present the attainable treatment and resource combinations before a preference is selected. Extensions to other reactor arrangements, attached-growth units, and resource-recovery connections can retain the same principle of defining component states and boundary streams before physical enforcement is imposed.

More extensive search comparisons can evaluate several initial controls, hybrid surrogate and mechanistic searches, and the fifteen-start basin-sensitivity design described in Chapter~\ref{ch:plant}. Recording state generation, fitting, prediction, projection, search, and verification costs would establish when surrogate development becomes worthwhile for repeated decisions. The central assessment would remain the treatment outcome and common-reference objective obtained for a given computational budget.

\section{Concluding perspective}

Physically constrained machine learning can connect rapid activated sludge prediction with the material accounting needed for process optimization. The reactor benchmark demonstrated consistent numerical enforcement across diverse learners. The structured regression added an inspectable response with favorable performance at the smallest data budget. The connected plant formulation brought unit consistency, solids storage, and independently evaluated operating choices into the same framework.

The dissertation's integrated contribution is a defined way to use these capabilities together. Statistical models provide rapid approximations, process knowledge supplies physical constraints, and mechanistic evaluation supplies the treatment response at selected controls. Future development can strengthen each stage through focused sampling, uncertainty assessment, dynamic balances, and plant measurements. This progression supports activated sludge optimization in which computational efficiency serves explicit treatment and resource objectives.
